
\documentclass{beamer}
\mode<presentation>
\usepackage{amsmath}
\usepackage{amssymb}
%\usepackage{advdate}
\usepackage{graphicx}
\graphicspath{{./figs/}}
\usepackage{adjustbox}
\usepackage{subcaption}
\usepackage{enumitem}
\usepackage{multicol}
\usepackage{mathtools}
\usepackage{listings}
\usepackage{url}
\def\UrlBreaks{\do\/\do-}
\usetheme{Boadilla}
\usecolortheme{lily}
\setbeamertemplate{footline}
{
  \leavevmode%
  \hbox{%
  \begin{beamercolorbox}[wd=\paperwidth,ht=2.25ex,dp=1ex,right]{author in head/foot}%
    \insertframenumber{} / \inserttotalframenumber\hspace*{2ex} 
  \end{beamercolorbox}}%
  \vskip0pt%
}
\setbeamertemplate{navigation symbols}{}

\providecommand{\nCr}[2]{\,^{#1}C_{#2}} % nCr
\providecommand{\nPr}[2]{\,^{#1}P_{#2}} % nPr
\providecommand{\mbf}{\mathbf}
\providecommand{\pr}[1]{\ensuremath{\Pr\left(#1\right)}}
\providecommand{\qfunc}[1]{\ensuremath{Q\left(#1\right)}}
\providecommand{\sbrak}[1]{\ensuremath{{}\left[#1\right]}}
\providecommand{\lsbrak}[1]{\ensuremath{{}\left[#1\right.}}
\providecommand{\rsbrak}[1]{\ensuremath{{}\left.#1\right]}}
\providecommand{\brak}[1]{\ensuremath{\left(#1\right)}}
\providecommand{\lbrak}[1]{\ensuremath{\left(#1\right.}}
\providecommand{\rbrak}[1]{\ensuremath{\left.#1\right)}}
\providecommand{\cbrak}[1]{\ensuremath{\left\{#1\right\}}}
\providecommand{\lcbrak}[1]{\ensuremath{\left\{#1\right.}}
\providecommand{\rcbrak}[1]{\ensuremath{\left.#1\right\}}}
\theoremstyle{remark}
\newtheorem{rem}{Remark}
\newcommand{\sgn}{\mathop{\mathrm{sgn}}}
\providecommand{\abs}[1]{\left\vert#1\right\vert}
\providecommand{\res}[1]{\Res\displaylimits_{#1}} 
\providecommand{\norm}[1]{\lVert#1\rVert}
\providecommand{\mtx}[1]{\mathbf{#1}}
\providecommand{\mean}[1]{E\left[ #1 \right]}
\providecommand{\fourier}{\overset{\mathcal{F}}{ \rightleftharpoons}}
%\providecommand{\hilbert}{\overset{\mathcal{H}}{ \rightleftharpoons}}
\providecommand{\system}[1]{\overset{\mathcal{#1}}{ \longleftrightarrow}}
%\providecommand{\system}{\overset{\mathcal{H}}{ \longleftrightarrow}}
	%\newcommand{\solution}[2]{\textbf{Solution:}{#1}}
%\newcommand{\solution}{\noindent \textbf{Solution: }}
\providecommand{\dec}[2]{\ensuremath{\overset{#1}{\underset{#2}{\gtrless}}}}
\newcommand{\myvec}[1]{\ensuremath{\begin{pmatrix}#1\end{pmatrix}}}
\let\vec\mathbf

\lstset{
%language=C,
frame=single, 
breaklines=true,
columns=fullflexible
}

\numberwithin{equation}{section}
\title{1.8.9}
\author{AI25BTECH11033 - Spoorthi N}
% \maketitle
% \newpage
% \bigskip
\begin{document}
{\let\newpage\relax\maketitle}
\renewcommand{\thefigure}{\theenumi}
\renewcommand{\thetable}{\theenumi}

\textbf{Question}

\noindent Find the area of the triangle with vertices $\Vec{A}(-1 , 0 , -2)$ , $\Vec{B}(0 , 2 , 1)$ and $\Vec{C}(-1 , 4 , 1)$. \,\,\,\,\, (12, 2023)


\textbf{solution}:\\


We are asked to find the area of the triangle with vertices:
\begin{align} 
\Vec{A}(-1 , 0 ,-2) \,\,\,\,  \Vec{B}(0,2,1) \,\,\,\, \Vec{C}(-1,4,1)
\end{align}

From the vectors;

\begin{align}
\overrightarrow{\vec{AB}}= \vec{B} - \vec{A} = (0-(-1), 2-0,1-(-2))=(1 , 2, 3)
\end{align}


\begin{align}
\overrightarrow{\vec{AC}}= \vec{C} - \vec{A}=(-1-(-1),,4-0,,1-(-2))=(0 , 4 , 3)
\end{align}

\begin{align}
\overrightarrow{\Vec{AB}}=(-1, 2 ,3) 
\end{align}

\begin{align}
\overrightarrow{\Vec{AC}}=(0,4,3) 
\end{align}

Write cross producet in matrix form:

\begin{align}
\vec{AB} \times \vec{AC} =
\begin{vmatrix}
\mathbf{i} & \mathbf{j} & \mathbf{k} \\
1 & 2 & 3 \\
0 & 4 & 3
\end{vmatrix}
\end{align}

Expand the determinant;
\begin{align}
= i(2 \cdot 3 - 3 \cdot 4) - j(1 \cdot 3 - 3 \cdot 0) + k(1 \cdot 4 - 2 \cdot 0)
\end{align}

\begin{align}
= i(6 - 12) - j(3 - 0) + k(4 - 0)
\end{align}

\begin{align}
= -6i - 3j + 4k
\end{align}
\begin{align}
\vec{AB} \times \vec{AC} = (-6, -3, 4)
\end{align}

Magnitude of cross product;
\begin{align}
|\vec{AB} \times \vec{AC}| = \sqrt{(-6)^2 + (-3)^2 + 4^2} 
= \sqrt{36 + 9 + 16} 
= \sqrt{61}
\end{align}

Area of the triangle;
\begin{align}
\text{Area} = \tfrac{1}{2} |\vec{AB} \times \vec{AC}|
= \tfrac{1}{2} \sqrt{61}
\end{align}
Area of the triangle

\begin{align}
\text{Area} = \tfrac{1}{2} \left| \vec{AB} \times \vec{AC} \right|
= \tfrac{1}{2} \sqrt{61}
\end{align}

\begin{align}
\frac{\sqrt{61}}{2} \approx 3.91
\end{align}

\end{document}
